\documentclass[12pt, french, latin.medieval, a4paper]{article}
\author{MBE MBE MINDJANA LOIC HENRI}
\title{GNPS et MDByTE}

\begin{document}
    Il existe plusieurs types d'équations matricielles importantes en mathématiques et en sciences de l'ingénieur. Voici une liste de certains types d'équations matricielles courants où $A$, $B$, $Q$, $C$ sont des matrices données et $X$ est une matrice inconnue

    \begin{itemize}
        \item Équation matricielle linéaire : 
            \begin{equation}
                AX = B
            \end{equation}
        \item Équation matricielle de Lyapunov :
            \begin{equation}
                AX + XA^T + Q = 0
            \end{equation}
        \item Équation matricielle de Riccati :
            \begin{equation}
                XA + AX - XBR^{-1}B^TX + Q = 0
            \end{equation}
        \item Équation matricielle d'algèbre linéaire :
            \begin{equation}
                AX = XB
            \end{equation}
        \item Équation matricielle de Sylvester :
            \begin{equation}
                AX + XB = C
            \end{equation}
    
        \item Équation matricielle de Stein : $(A + XBX)^2 - (A^2 + XB^2X) = 0$
            \begin{equation}
                AX + XB = C
            \end{equation}
        \item Équation matricielle de DARE (Discrete-time Algebraic Riccati Equation) :
            \begin{equation}
                X = A^TXA - (A^TXB)(R + B^TXB)^{-1}(B^TXA) + Q
            \end{equation}
        \item Équation matricielle de CARE (Continuous-time Algebraic Riccati Equation) : 
            \begin{equation}
                AX + XA^T - XBR^{-1}B^TX + Q = 0
            \end{equation}
        \item Équation matricielle de Fokker-Planck :
            \begin{equation}
                \frac{\partial X}{\partial t} = D \frac{\partial^2 X}{\partial x^2} - \frac{\partial}{\partial x}(AX) + BX
            \end{equation}
        \item Équation matricielle de diffusion : $\nabla^2$ étant l'opérateur de Laplace-Beltrami
            \begin{equation}
                \frac{\partial X}{\partial t} = D \nabla^2 X
            \end{equation}
        
    \end{itemize}

\end{document}